\begin{scanpage}{0182}{164}
The coordinate system was selected in such a way that its origin would lie near the body, and the source Q would be located in the plane $x=0$. If the distance between the body and the source is a great deal larger than the body's dimensions, then the incident wave in the vicinity of the body may be investigated as a plane wave. Let us represent it in the form
\[
E_x=E_{0x}e^{ik(y\sin\gamma+z\cos\gamma)},\quad E_y=0.\tag{26.01}
\]
Here $\gamma$ is the angle between the normal $\mathbf N$ to the wave front and the $z$ axis.

Now let us place in front of the source, parallel to the radiated wave front, a polarizer P which transformed linear polarized radiation into a circularly polarized wave. Let the wave passing through the polarizer with an electric vector $E_\eta$ lag in phase by $90^\circ$ behind the wave with an electric vector $E_\tau$ (Figure 70). In this case, the polarizer achieves a clockwise rotation$^{(1)}$. As a result, the incident wave field at the coordinate origin will equal
\[
E_x=\frac{e^{i\pi/4}}{\sqrt2}E_{0x},\quad H_x=\frac{e^{-i\pi/4}}{\sqrt2}E_{0x}.\tag{26.02}
\]
The field scattered by the body may be represented in the wave zone in the following way:
\[
\left.\begin{aligned}E_\varphi=-H_\vartheta&=\frac{iaE_{0x}}2\frac{e^{i\pi/4}}{\sqrt2}\overline\Sigma(\vartheta,\varphi)\frac{e^{ikR}}R,\\
E_\vartheta=H_\varphi&=\frac{iaE_{0x}}2\frac{e^{-i\pi/4}}{\sqrt2}\Sigma(\vartheta,\varphi)\frac{e^{ikR}}R.
\end{aligned}\right\}\tag{26.03}
\]
where $a$ is a certain length characterizing the body's size and $\overline\Sigma(\vartheta,\varphi)$ and $\Sigma(\vartheta,\varphi)$ are unknown angular functions. In the general case, the
\par\bigskip\hrule\smallskip
$(1)$Footnote appears on page 174.
\end{scanpage}
